\section{The multiplicity problem}
\label{sec:introduction}

For a nonzero convergent power series
$f\in\C\{z_1,\ldots,z_N\}$ with $f(0)=0$, the order $\ord_0f$ is
the least total degree of a monomial with nonzero coefficient. When $f$
is reduced, this integer is the multiplicity of the hypersurface germ
$V(f)=\{f=0\}$. Two such germs are \emph{ambiently homeomorphic} if a
homeomorphism between neighborhoods of the origin in $\C^N$ fixes the
origin and carries one hypersurface to the other. Zariski's multiplicity
question asks whether this embedded topological equivalence determines
multiplicity~\cite{Zariski1971}. It asks, in particular, whether topology
can distinguish a hypersurface with a nonzero quadratic term from one
whose first nonzero terms are cubic.

For an irreducible plane curve, the classical relation between embedded
topology and the characteristic data of a branch determines its
multiplicity~\cite{Zariski1932}. Summing the branch multiplicities gives
the same conclusion for a reduced plane curve. In every ambient dimension,
the A'Campo--L\^e theorem shows that a reduced hypersurface germ ambiently
homeomorphic to a smooth one is itself smooth~\cite{ACampo1973,Le1973};
see~\cite[Theorem~2.2]{Sampaio2025Milnor}. Thus multiplicity one cannot
occur in an unequal-multiplicity pair. Eyral surveys the multiplicity
problem and its further positive cases~\cite{Eyral2007}.
In higher dimensions, the Milnor
fibration gives powerful topological invariants of an isolated
singularity~\cite{Milnor1968}. Its integral middle homology carries both
a monodromy operator and a Seifert form. The latter records linking of
cycles in neighboring fibers and, in a suitable high-dimensional range,
classifies the embedded link. Multiplicity, by contrast, is read directly
from the lowest-degree terms of an equation. Our construction separates
these two kinds of information.

A polynomial is \emph{weighted homogeneous with positive rational
weights} if there are positive integers $a_1,\ldots,a_N,d$ such that
$f(t^{a_1}z_1,\ldots,t^{a_N}z_N)=t^d f(z)$; the normalized weights are
$w_i=a_i/d$. We prove the following existence statement.

\begin{theorem}\label{thm:main}
There is a finite integer $N>3$, divisible by eight, and there are reduced
polynomials $f_1,f_2\in\R[z_1,\ldots,z_N]$, weighted homogeneous for
positive rational weights, such that:
\begin{enumerate}[label=\textup{(\roman*)}]
\item each $f_j$ has an isolated critical point at the origin;
\item $\ord_0f_1=2$ and $\ord_0f_2=3$;
\item there is a homeomorphism germ of ambient pairs
\[
 (\C^N,V(f_1),0)\ \cong\ (\C^N,V(f_2),0).
\]
\end{enumerate}
\end{theorem}

Theorem~\ref{thm:main} gives a negative answer to the embedded Zariski
multiplicity conjecture. A single finite ambient dimension suffices for
this purpose. The construction determines such a dimension by exact
integer arithmetic; it is not intended to minimize it.

The same examples settle the corresponding questions for topological right
equivalence. Holomorphic function germs $f,g:(\C^N,0)\to(\C,0)$ are
\emph{topologically right equivalent} if $f=g\circ\varphi$ for a
homeomorphism germ $\varphi:(\C^N,0)\to(\C^N,0)$. The right-topological
multiplicity conjecture asks whether reduced germs related in this way have
equal multiplicities; its mod-two version asks only for equal multiplicity
parity \cite[Conjecture~2 and Conjecture~Z mod~2]{Sampaio2025Milnor}.
Corollary~\ref{topo:right-equivalence} gives negative answers to both
questions for the pair in Theorem~\ref{thm:main}. The source homeomorphism
there may differ from the ambient map constructed below and is not claimed
to be holomorphic or bi-Lipschitz.

A further consequence concerns Arnold's corank problem. For a holomorphic
germ $f:(\C^N,0)\to(\C,0)$ with a critical point at the origin, write
\[
 \operatorname{corank}_0(f)
       =N-\rank_{\C}\operatorname{Hess}_0(f),
\]
where $\operatorname{Hess}_0(f)$ is the Hessian matrix at the origin.
Arnold's question asks whether this number is preserved when a
homeomorphism germ of $(\C^N,0)$ carries $V(f)$ to $V(g)$ for reduced
holomorphic germs $f,g$ with critical points at the origin
\cite[Problem~1]{Sampaio2026}. Corollary~\ref{topo:arnold-corank} gives
a negative answer, even for the isolated weighted homogeneous germs
with real coefficients in Theorem~\ref{thm:main}.

The examples also answer a question about the lowest-degree part of a
function. For a nonzero holomorphic germ $g$ vanishing at the origin,
let $g_*$ denote its lowest nonzero homogeneous component for ordinary
total degree. Its \emph{initial-form Milnor fibre} has the homotopy type
of the global unit fibre $g_*^{-1}(1)$; the full polynomial $g_*$ is
retained, even when it is nonreduced
\cite[Proposition~1.6 in arXiv v3]{BudurEtAl2022ContactLoci}.
Motivated by the relationship between contact loci and monodromy,
Budur, Fern\'andez de Bobadilla, L\^e, and Nguyen conjectured that these
fibres are homotopy equivalent for embedded topologically equivalent
function germs~\cite[Conjecture~1.7 in arXiv v3]{BudurEtAl2022ContactLoci}.
Sampaio studies the explicit right-equivalence formulation
\cite[Conjecture~4]{Sampaio2025Milnor}.
Corollary~\ref{topo:tangent-milnor} disproves both formulations after adding one common
square: the germs $f_j(z)+t^2$ in $\C^{N+1}$ remain topologically right
equivalent and have isolated critical points, but their initial-form Milnor fibres are
connected and have two connected components, respectively. Even their
zeroth integral homology groups are therefore different.

\subsection*{Families, metric equivalence, and the present question}

The question for an arbitrary pair must be distinguished from
multiplicity in a family. L\^e and Ramanujam showed that a one-parameter holomorphic deformation
of isolated hypersurface singularities with constant Milnor number is
topologically constant when the ambient complex dimension differs from
three~\cite{LeRamanujam1976}. In the excluded case $\C^3$, the companion
\cite[Theorem~1.1]{OpenAIMuConstant2026} proves, after shrinking,
topological right-triviality for holomorphic one-parameter families $f_t$ with
$f_t(0)=0$, an isolated critical point at the origin for every $t$,
and constant $\mu(f_t,0)$.
Greuel and O'Shea proved equimultiplicity
for the corresponding deformation problem at an isolated quasihomogeneous
singularity~\cite{Greuel1986,OShea1987}. Fern\'andez de Bobadilla and
Pe\l ka subsequently established equimultiplicity for continuous
families with finite constant Milnor number, including coefficientwise
continuous families of formal power series
\cite[Theorem~1.1 and Corollary~1.3]{BobadillaPelka2022}.
Their Example~8.6 gives topologically equivalent germs in $\C^4$, both
of multiplicity two, that cannot be joined by a continuous
$\mu$-constant family. They obtain the embedded equivalence by adding
two independent squares to plane-curve germs with the same integral
Seifert form, a precedent for the passage from integral forms to
embedded topology used here.
Thus an arbitrary equivalent pair does not
reduce to the family problem. The two different multiplicities in
Theorem~\ref{thm:main} therefore preclude a connecting $\mu$-constant
family.

There are also positive results for pairs under additional geometric
hypotheses. Yau established multiplicity invariance for pairs of isolated
quasihomogeneous surface singularities in $\C^3$
\cite[Theorem~C]{Yau1988}, crediting the detailed proofs jointly with
Xu in the article's added note. Metric and low-dimensional hypotheses
give further positive results. Fernandes, Jelonek, and Sampaio proved that multiplicity two is preserved
by ambient bi-Lipschitz equivalence, and that Hessian corank is preserved
by ambient homeomorphisms in three complex variables
\cite[Corollary~3.5 and Theorem~4.4]{FernandesJelonekSampaio2025}.
The homeomorphism in Theorem~\ref{thm:main} is topological and its
ambient dimension is larger than three. Sampaio's more recent
corank-parity theorem gives an ambient topological invariant in every
dimension~\cite[Theorem~4.1]{Sampaio2026}.
Corollary~\ref{topo:arnold-corank} is consistent with this theorem:
the two Hessian ranks are even, so the unequal coranks have the same parity.
These comparisons describe the precise topological setting of the
counterexample.

\subsection*{The main ideas}

The main difficulty is to preserve different ordinary orders while
making the \emph{integral Seifert forms} coincide. We compare the
spectra, finite multisets of rational exponents whose exponentials give
the monodromy eigenvalues. Equality of spectra alone does not accomplish
the desired integral identification. For the isolated weighted homogeneous
singularities considered here, in a fixed ambient dimension, the spectrum
determines the real Seifert form, through the work of Steenbrink and
N\'emethi
\cite{Steenbrink1977,Steenbrink1989,Nemethi1995,BalnojanHertling2019}.
For sums of chain singularities, the integral-monodromy results of
Hertling and Mase, building on Orlik and Randell, separately determine
the integral module with its monodromy operator
\cite{Orlik1972,OrlikRandell1977,HertlingMase2022}.
The two resulting identifications need not be compatible. The arithmetic
part of the proof produces a single integral map preserving the whole
Seifert form.

\paragraph{1. Seeds from a finite relation.}
We begin with two multiplicative identities among explicitly specified
positive rational numbers. A finite valuation matrix certifies an
integer relation whose coefficient sum is positive. A decreasing
integer-remainder construction completes each required weight to an
isolated chain polynomial. The two collections of chains have different
numbers of variables; balancing the difference by squares on one side
produces multiplicity two there, while common cubes give multiplicity
three on the other. The product identities and a triple-angle identity
show that the spectra agree modulo two. The two classical classification
results just described then supply a real Seifert-form isometry and an
integral monodromy-module isomorphism
(Section~\ref{seed:section}).

\paragraph{2. Compatible local and rational forms.}
After adding common fifth and seventh powers, the monodromy has finite
odd order and no eigenvalue one. In even ambient dimension, this permits
an integral symmetrization of each Seifert form that retains the
monodromy and recovers the original form. We prove that a first tensor
square of these symmetric data has equivariant isometries over every
$\Z_\ell$. At odd primes, square roots and sums of two squares in a
finite operator algebra produce the maps. At two, the operator splits
the lattice into dual pairs and unramified hermitian modules. A
quadratic-form local--global argument then supplies an equivariant
rational isometry (Section~\ref{local:section}).

\paragraph{3. A single integral isometry.}
The rational map and the local maps can still differ by nontrivial
eigenspace determinants. Real coefficients provide integral involutions
reversing monodromy. A second tensor square pairs each map with its
conjugate by these involutions; an exact determinant calculation cancels
the discrepancies on every eigenspace. Six common cubic variables make
every relevant hermitian form indefinite. Two further common pure
powers, of degrees eleven and thirteen, eliminate real eigenvalues.
The determinant-one isometry group is now a product of special unitary
groups. The strong-approximation method of Eichler, Kneser, and Platonov
\cite{Platonov1969,Platonov1970,Rapinchuk2014} gives a rational group
element meeting all prescribed local lattice conditions, and hence an
integral Seifert-form isometry (Section~\ref{global:section}).

\paragraph{4. From forms to ambient germs.}
The Thom--Sebastiani theorem and Sakamoto's integral Seifert-form formula
realize the tensor operations by polynomial sums in disjoint variables
\cite{SebastianiThom1971,Sakamoto1974}. The last two pure powers also
force $\det(1-T)=1$ for the final monodromy $T$. The links are thus
homotopy spheres, and the Milnor fibration verifies that they are simple
knots. Levine's classification converts the integral Seifert congruence
into embedded link equivalence~\cite{Levine1970}. Weighted radial
coordinates extend that equivalence to a homeomorphism of ambient germs
(Section~\ref{topo:section}). Figure~\ref{fig:proof-route} collects
these stages and their distinct outputs.

Section~\ref{topo:section} then applies the results of King and Saeki
to obtain topological right equivalence. Real coefficients remove the
complex-conjugation alternative in that transfer
\cite{King1978,Saeki1989,Sampaio2025Milnor}.
Adding a square in one new variable preserves right equivalence and
gives the different connectedness of the initial-form Milnor fibres.

\begin{figure}[htbp]
\centering
\begingroup
\setlength{\tabcolsep}{0pt}
\renewcommand{\arraystretch}{1.08}
\begin{tabular}{c}
\fbox{\parbox{.88\linewidth}{\centering
\textbf{Arithmetic relations and chain polynomials}\\
Equal spectra modulo $2$; ordinary orders $2$ and $3$.}}\\[3pt]
$\Downarrow$\\[3pt]
\fbox{\parbox{.88\linewidth}{\centering
\textbf{Two separate seed identifications}\\
Real Seifert forms \quad and \quad integral monodromy modules.}}\\[3pt]
$\Downarrow\quad$ common fifth and seventh powers;\\
\textit{symmetrize and take one tensor square}\\[3pt]
\fbox{\parbox{.88\linewidth}{\centering
\textbf{Equivariant isometries over $\Q$ and every $\Z_\ell$}\\
A second tensor square cancels their eigenspace determinant defects.}}\\[3pt]
$\Downarrow\quad$ common cubic, eleventh, and thirteenth powers\\[3pt]
\fbox{\parbox{.88\linewidth}{\centering
\textbf{One integral Seifert-form isometry}\\
Indefinite special unitary groups and strong approximation.}}\\[3pt]
$\Downarrow\quad$ simple spherical knots and weighted radial coordinates\\[3pt]
\fbox{\parbox{.88\linewidth}{\centering
\textbf{Ambiently homeomorphic hypersurface germs}\\
Disjoint sums preserve the two different ordinary orders.}}
\end{tabular}
\endgroup
\caption{The proof keeps analytic order separate from the successive
identifications of topological data. The two seed identifications are
independent; the local and global arithmetic makes them compatible.}
\label{fig:proof-route}
\end{figure}


Section~\ref{sec:background} fixes the conventions and exact classical
inputs. Appendix~\ref{app:certificate} contains the complete integer
program for the finite rank certificate; a separately supplied exact
relation can be verified by trial division alone. All construction,
local algebra, determinant calculations, and the passage from local
lattices to an integral map are proved in the text. The classical
singularity-theoretic, arithmetic, and topological classification
theorems are invoked after checking their hypotheses.

\subsection*{Conventions}

Every sum combining singularities uses disjoint sets of variables.
Superscript $[i]$ marks a fresh copy of a polynomial. An isometry of
lattices with operators is required to intertwine those operators.
The transpose of a matrix is $M^{\mathsf t}$. A lattice is a finite-rank
free abelian group; a bilinear form is unimodular if its determinant is
$\pm1$.
